\chapter{Stage 5: The Nine-Signal Verifier}
\label{ch:verifier}

A tier has produced a fresh answer. Stage 5 is where the system decides how much
to believe it. This is the heart of \caem{}'s trust machinery and the direct
answer to the overconfidence flaw of \Cref{ch:problem}: rather than read the
model's own confidence, which we saw is barely better than chance, the verifier
gathers many independent signals about the answer and weighs them together. This
chapter covers how each raw signal is produced. The next chapter, \Cref{ch:composite},
covers how they are calibrated, fused into one number, and turned into a decision.

\begin{intuition}[A panel of specialists, not a single judge]
No single test reliably catches a confident falsehood. So the verifier convenes a
panel. One specialist asks whether the model was internally sure. Another asks
whether the model agrees with itself when made to answer several times. A third
checks the answer against retrieved evidence, sentence by sentence. A fourth asks
whether the answer even addresses the question. Each specialist is fallible alone,
but a confident lie has to fool all of them at once, and that is much harder than
fooling any one. The verifier's strength is not any single signal; it is the
diversity of the panel.
\end{intuition}

\section{Four families, nine headline signals}
\label{sec:verifier-families}

The signals group into four families, each probing a different way an answer can be
wrong. \Cref{fig:families} is the map to keep in view for the whole chapter.

\begin{figure}[t]
\centering
\begin{tikzpicture}[
  font=\small\sffamily,
  fam/.style={rectangle, rounded corners=5pt, line width=1.1pt, align=left,
    inner sep=6pt, text width=5.2cm, minimum height=2.0cm,
    drop shadow={shadow xshift=0.3ex, shadow yshift=-0.3ex, fill=black!10}},
  core/.style={rectangle, rounded corners=4pt, draw=cbInk!55, fill=cbInk!8,
    line width=1pt, align=center, inner sep=5pt, text width=2.6cm, font=\footnotesize},
  flow/.style={-{Stealth[length=2.4mm]}, line width=1pt, draw=cbInk!65},
]
  \node[fam, draw=cbIntuition!75!black, fill=cbIntuition!8] (f1) at (0,2.1)
    {\textbf{\textcolor{cbIntuition!75!black}{Internal calibration}}\\[2pt]
     {\footnotesize Was the model itself sure?}\\[3pt]
     {\footnotesize $\utoken$, $\udrop$, $\uint$}};
  \node[fam, draw=cbScenario!70!black, fill=cbScenario!8] (f2) at (6.4,2.1)
    {\textbf{\textcolor{cbScenario!60!black}{Self-consistency}}\\[2pt]
     {\footnotesize Does it agree with itself?}\\[3pt]
     {\footnotesize $\savg$, $\pentail$ \;\;(over $M{=}3$ chains)}};
  \node[fam, draw=cbExample!70!black, fill=cbExample!10] (f3) at (0,-0.9)
    {\textbf{\textcolor{cbExample!50!black}{External grounding}}\\[2pt]
     {\footnotesize Does the evidence support it?}\\[3pt]
     {\footnotesize $\pgmax$, $\pgmean$, $\pgatom$}};
  \node[fam, draw=cbNumbers!75!black, fill=cbNumbers!10] (f4) at (6.4,-0.9)
    {\textbf{\textcolor{cbNumbers!65!black}{Relevance}}\\[2pt]
     {\footnotesize Does it answer the question?}\\[3pt]
     {\footnotesize $\qarel$}};
  \node[core] (comp) at (3.2,-3.0) {to the calibrated composite (\Cref{ch:composite})};
  \foreach \f in {f1,f2,f3,f4} \draw[flow] (\f) -- (comp);
\end{tikzpicture}
\caption[The four signal families]{The nine headline signals, grouped into four
families, each probing a different failure mode. Every family feeds the calibrated
composite of \Cref{ch:composite}. A confident wrong answer must satisfy all four at
once to be trusted. (Schematic.)}
\label{fig:families}
\end{figure}

\begin{pitfall}[``Nine signals'' is exact: the composite computes more than it uses]
The verifier is named for the nine signals of \Cref{fig:families}, and that is
exactly how many drive a decision: the calibrated composite of \Cref{ch:composite}
reads precisely those nine. The verifier nonetheless \emph{computes} more than it
uses. Two entity checks, an alias overlap and a head-noun consistency
(\Cref{sec:verifier-entity}), were trialed in an earlier revision and then retired
from the composite; two original signals, a semantic-entropy measure and a
contradiction probability (\Cref{sec:verifier-retired}), were retired as well. All
four are still computed and recorded for diagnostics, but none contributes to the
trust score in the deployed configuration. So ``nine'' is not a headline
simplification; it is the exact count of signals that move the gate.
\end{pitfall}

\section{How the signals are produced, cheaply}
\label{sec:verifier-dataflow}

Several signals would be expensive computed naively, so the verifier produces them
from a small number of shared generations. The key economy is that the model is
made to answer the same query a few times, and those few generations are reused
across multiple signals. \Cref{fig:verifier-flow} traces the production.

\begin{figure}[t]
\centering
\resizebox{\textwidth}{!}{%
\begin{tikzpicture}[
  font=\small\sffamily,
  inp/.style={rectangle, rounded corners=2pt, draw=cbInk!55, fill=cbInk!6,
    line width=1pt, align=center, inner sep=5pt, text width=2.3cm},
  gen/.style={rectangle, rounded corners=4pt, draw=cbScenario!70!black, fill=cbScenario!12,
    line width=1pt, align=center, inner sep=5pt, text width=3.0cm,
    drop shadow={shadow xshift=0.3ex, shadow yshift=-0.3ex, fill=black!10}},
  sigb/.style={rectangle, rounded corners=3pt, draw=cbFormula!70!black, fill=cbFormula!10,
    line width=0.9pt, align=center, inner sep=3pt, text width=2.0cm, font=\footnotesize,
    minimum height=0.7cm},
  flow/.style={-{Stealth[length=2.4mm]}, line width=1pt, draw=cbInk!70},
]
  \node[inp] (qa) at (0,0) {query + fresh answer};

  \node[gen] (mchain) at (3.8,2.4) {\textbf{$M{=}3$ chains}\\{\footnotesize one generate call, $T{=}0.7$}};
  \node[gen] (drop)   at (3.8,0.0) {\textbf{$K{=}5$ dropout passes}\\{\footnotesize plus teacher forcing}};
  \node[gen] (ret)    at (3.8,-2.4) {\textbf{retrieve 20 $\to$ rerank 3}\\{\footnotesize top passages}};

  \node[sigb, fill=cbScenario!12, draw=cbScenario!70!black] (savg)  at (8.0,3.2) {$\savg$};
  \node[sigb, fill=cbScenario!12, draw=cbScenario!70!black] (pent)  at (8.0,2.4) {$\pentail$};
  \node[sigb, fill=cbScenario!12, draw=cbScenario!70!black] (ehead) at (8.0,1.6) {entity-head};
  \node[sigb, fill=cbIntuition!12, draw=cbIntuition!75!black] (utok) at (8.0,0.5) {$\utoken$};
  \node[sigb, fill=cbIntuition!12, draw=cbIntuition!75!black] (udro) at (8.0,-0.3) {$\udrop\!\to\!\uint$};
  \node[sigb, fill=cbExample!12, draw=cbExample!70!black] (pg)   at (8.0,-1.6) {$\pgmax$\\$\pgmean$\\$\pgatom$};
  \node[sigb, fill=cbNumbers!12, draw=cbNumbers!75!black] (qar)  at (8.0,-2.4) {$\qarel$};
  \node[sigb, fill=cbExample!12, draw=cbExample!70!black] (alias) at (8.0,-3.2) {alias-overlap};

  \node[inp, text width=2.2cm] (out) at (11.6,0) {nine composite signals\\{\footnotesize(+ recorded diagnostics)}};

  \draw[flow] (qa) -- (mchain);
  \draw[flow] (qa) -- (drop);
  \draw[flow] (qa) -- (ret);
  \draw[flow] (mchain) -- (savg);
  \draw[flow] (mchain) -- (pent);
  \draw[flow] (mchain) -- (ehead);
  \draw[flow] (drop) -- (utok);
  \draw[flow] (drop) -- (udro);
  \draw[flow] (ret) -- (pg);
  \draw[flow] (ret) -- (qar);
  \draw[flow] (ret) -- (alias);
  \foreach \s in {savg,pent,ehead,utok,udro,pg,qar,alias} \draw[flow] (\s) -- (out);
\end{tikzpicture}}
\caption[The verifier data flow]{The verifier produces all signals from three
shared computations: one generation of $M{=}3$ reasoning chains (reused by three
signals), $K{=}5$ dropout passes plus a teacher-forcing pass for the internal
signals, and one retrieve-and-rerank for the grounding and relevance signals.
Generating the three chains once and reusing them is the central economy of the
stage. (Schematic.)}
\label{fig:verifier-flow}
\end{figure}

\begin{precise}[The canonical answer claim: one surface for every signal]
Before any signal reads the answer, the answer is rewritten into a \emph{canonical
claim}, and every signal that touches the answer reads that same canonical form.
A multiple-choice letter such as ``C'' is expanded to its full option text; a bare
entity such as ``Paris'' is wrapped as ``Answer: Paris.''; an already-declarative
answer such as a fact-check verdict passes through unchanged. The reason is
uniformity: the grounding signals, the entailment signal, and the relevance signal
must all score the \emph{same} propositional surface, or their numbers would not be
comparable across benchmarks and the single calibrated threshold of
\Cref{ch:composite} could not govern them all. Canonicalisation is the small,
unglamorous step that makes the whole panel speak one language.
\end{precise}

\section{Family one: internal calibration}
\label{sec:verifier-internal}

The first family reads the generator itself, at no external cost. It contains three
signals. The token confidence $\utoken$ is the geometric mean of the per-token
probabilities of the answer, the same quantity we met in \Cref{ch:upre}, here
computed by teacher-forcing the produced answer rather than from a short decode. The
dropout disagreement $\udrop$ measures stability under perturbation: the model is
made to answer five times with dropout left switched on, and the signal is the
fraction of those five samples that disagree with the plurality answer. The two are
combined into a single internal confidence,
\[
  \uint \;=\; 0.5\,\utoken \;+\; 0.5\,(1 - \udrop),
\]
weighting token reliability and answer stability equally.

\begin{lstlisting}[caption={Dropout disagreement, condensed from \texttt{caem/verification/verifier.py}.}, label={lst:udropout}]
samples = model.generate(num_return_sequences=5, do_sample=True,    # K=5, dropout ON
                         temperature=0.7)                           # model.train() mode
answers = [extract_answer(s) for s in samples]                      # short form of each
plurality = most_common(answers)
u_dropout = 1.0 - answers.count(plurality) / 5                      # fraction that disagree
\end{lstlisting}

\begin{intuition}[Two ways to be unsure]
The two internal signals catch different doubts. Token confidence is low when the
model hedged \emph{within} a single answer, spreading probability over many next
tokens. Dropout disagreement is high when the model is unstable \emph{across}
answers, landing on different conclusions when nudged. A model can be fluent within
one answer yet unstable across attempts, or stable across attempts yet hesitant
within each; combining the two catches both. Both are cheap and require no external
model, which is why they anchor the panel.
\end{intuition}

\section{Family two: self-consistency}
\label{sec:verifier-consistency}

The second family makes the model talk to itself. It generates $M=3$ reasoning
chains for the query, sampled at a temperature of $0.7$, in a single batched call,
and two signals read those three chains. The semantic consistency $\savg$ is the
mean similarity across all unique pairs of chains, measured by embedding each chain
and taking cosine similarity; high values mean the model keeps reaching the same
place by different routes~\cite{wang2023selfconsistency}. The chain-to-answer
entailment $\pentail$ asks a sharper question: does each chain actually
\emph{entail} the served answer? Each chain is run through the natural-language-
inference judge against the canonical answer, and the signal is the mean entailment.

\begin{intuition}[Agreement is not the same as support]
Why two signals from the same three chains? Because chains can agree with each other
yet not support the answer that was served. Imagine three chains that all reason
toward one conclusion, but the answer the system actually returned was a different
one: $\savg$ would be high (the chains agree) while $\pentail$ would be low (they do
not entail the served answer). That gap is a real failure mode, an answer
disconnected from the reasoning meant to justify it, and it is exactly what
$\pentail$ exists to catch. Generating the chains once and reading them two ways is
both cheaper and more informative than two separate tests.
\end{intuition}

\section{Family three: external grounding}
\label{sec:verifier-grounding}

The third family is the strongest evidence that an answer is true rather than merely
confident: does retrieved evidence support it? The verifier retrieves twenty
candidate passages for the query, reranks them with a cross-encoder to the best
three, and scores the canonical answer against those three passages with the
natural-language-inference judge. Three signals read the result, and they differ in
how forgiving they are.

\begin{figure}[t]
\centering
\begin{tikzpicture}[
  font=\small\sffamily,
  ans/.style={rectangle, rounded corners=4pt, draw=cbInk!55, fill=cbInk!6,
    line width=1pt, align=center, inner sep=5pt, text width=3.2cm},
  pass/.style={rectangle, rounded corners=3pt, draw=cbExample!60!black, fill=cbExample!8,
    line width=0.9pt, align=left, inner sep=4pt, text width=2.6cm, font=\footnotesize},
  res/.style={rectangle, rounded corners=3pt, draw=cbFormula!70!black, fill=cbFormula!10,
    line width=0.9pt, align=center, inner sep=4pt, text width=3.4cm, font=\footnotesize},
  atom/.style={rectangle, rounded corners=2pt, draw=cbNumbers!70!black, fill=cbNumbers!10,
    line width=0.8pt, align=left, inner sep=3pt, text width=3.0cm, font=\scriptsize},
]
  \node[ans] (a) at (0,1.0) {canonical answer claim};
  \node[pass] (p1) at (3.6,2.0) {passage 1\hfill\textbf{0.94}};
  \node[pass] (p2) at (3.6,1.0) {passage 2\hfill\textbf{0.55}};
  \node[pass] (p3) at (3.6,0.0) {passage 3\hfill\textbf{0.40}};
  \draw[-{Stealth[length=2mm]}, cbInk!55] (a) -- (p1);
  \draw[-{Stealth[length=2mm]}, cbInk!55] (a) -- (p2);
  \draw[-{Stealth[length=2mm]}, cbInk!55] (a) -- (p3);

  \node[res, fill=cbFormula!14] (gmax)  at (8.1,2.0) {$\pgmax = \max = \mathbf{0.94}$\\{\scriptsize best single passage}};
  \node[res] (gmean) at (8.1,1.0) {$\pgmean = \mathrm{mean} = \mathbf{0.63}$\\{\scriptsize average support}};
  \node[res] (gatom) at (8.1,-0.6) {$\pgatom = \min_{\text{atoms}} = \mathbf{0.40}$\\{\scriptsize weakest atomic fact}};
  \draw[-{Stealth[length=2mm]}, cbInk!40] (p1) -- (gmax);
  \draw[-{Stealth[length=2mm]}, cbInk!40] (p2) -- (gmean);

  \node[atom] (atoms) at (3.6,-1.7) {decompose answer:\\$\bullet$ atom A \hfill 0.91\\$\bullet$ atom B \hfill 0.40\\$\bullet$ atom C \hfill 0.78};
  \draw[-{Stealth[length=2mm]}, cbNumbers!60!black] (a) to[bend right=18] (atoms);
  \draw[-{Stealth[length=2mm]}, cbNumbers!60!black] (atoms) -- (gatom);
\end{tikzpicture}
\caption[The grounding trio]{Three readings of the same evidence, in rising
strictness. $\pgmax$ takes the single best passage and is the most forgiving;
$\pgmean$ averages the top three; $\pgatom$ decomposes the answer into atomic facts,
scores each against the passages, and reports the \emph{weakest} one. A long answer
with one unsupported fact passes the first two and fails the third. (Numbers are
illustrative.)}
\label{fig:grounding-trio}
\end{figure}

\begin{precise}[Max, mean, and the weakest link]
The maximum grounding $\pgmax$ is the entailment of the single best-supporting
passage, the most forgiving reading; it is also the signal the early-exit gate of
\Cref{sec:verifier-earlyexit} watches. The mean grounding $\pgmean$ averages over
the three passages, a steadier estimate of overall support. The atomic grounding
$\pgatom$ is the strict one: it decomposes the answer into atomic facts in the
FActScore style~\cite{min-etal-2023-factscore}, scores each fact against the
passages, and returns the \emph{weakest} fact's support, on the principle that one
unsupported claim in an otherwise-grounded answer is still a fabrication. Because
decomposition is meaningless for very short answers, it is gated by length: answers
under twelve tokens skip it and fall back to the mean. The three together span the
range from ``is there any support at all?'' to ``is every part supported?''
\end{precise}

\subsection*{The two natural-language-inference judges}

The grounding and entailment signals all rest on a judge that scores whether a
premise entails a hypothesis. \caem{} uses two judges and dispatches between them by
the length of the hypothesis, as \Cref{fig:nli-dispatch} shows.

\begin{figure}[t]
\centering
\begin{tikzpicture}[font=\small\sffamily]
  % token budget bar
  \fill[cbDefense!18] (0,0) rectangle (1.2,0.6);
  \node[font=\scriptsize] at (0.6,0.3) {4};
  \fill[cbScenario!20] (1.2,0) rectangle (4.0,0.6);
  \node[font=\scriptsize] at (2.6,0.3) {100 premise};
  \fill[cbIntuition!18] (4.0,0) rectangle (12.0,0.6);
  \node[font=\scriptsize] at (8.0,0.3) {408 hypothesis budget};
  \draw[cbInk!50, line width=0.8pt] (0,0) rectangle (12.0,0.6);
  \node[font=\scriptsize, anchor=west] at (0,0.95) {MiniCheck encoder window: $512 = 4 \text{ overhead} + 100 \text{ premise} + 408 \text{ hypothesis}$};
  % dispatch
  \node[font=\small\sffamily, anchor=west] at (0,-0.9)
    {hypothesis $\le 408$ tokens \;$\Rightarrow$\; \textcolor{cbScenario!60!black}{\textbf{MiniCheck}} (fast specialist)};
  \node[font=\small\sffamily, anchor=west] at (0,-1.6)
    {hypothesis $> 408$ tokens \;$\Rightarrow$\; \textcolor{cbIntuition!70!black}{\textbf{Frozen Qwen}} (long context, Platt-scaled to MiniCheck's scale)};
\end{tikzpicture}
\caption[The length-dispatched NLI judges]{Entailment is judged by a fast specialist
for short hypotheses and a long-context model for long ones. The split point of
$408$ tokens is the hypothesis budget left once the specialist's $512$-token window
reserves four tokens of overhead and a hundred tokens of premise. The long-context
judge is Platt-scaled on a shared fold so both judges report on one comparable
probability scale. In the deployed run, however, the long-context judge was set aside
because its post-hoc probability fit did not clear the required correlation floor, so
the shipped verifier ran the fast specialist for every prediction and truncated the
rare longer hypotheses at that checker's window; this is recorded among the threats. (Schematic.)}
\label{fig:nli-dispatch}
\end{figure}

Two further refinements keep the grounding honest. For refutation tasks such as
fact-checking, a \emph{directional} scorer rewrites the hypothesis so the judge is
asked whether the passages support the \emph{truth} of the claim rather than the
claim as stated, which is the correct question when the right answer may be
``refutes.'' And the ensemble of judges is combined conservatively: for entailment
the minimum across judges is taken, so every judge must agree before the answer is
counted as grounded.

\section{Family four: relevance, and two entity guards}
\label{sec:verifier-entity}

The fourth family contains a single headline signal, the question-answer relevance
$\qarel$. A cross-encoder scores whether the answer actually addresses the question,
turning its logit into a probability. This closes a subtle failure mode: an answer
that is fluent, internally consistent, and even grounded in a retrieved passage, but
simply off-topic, answering a different question than the one asked. Grounding alone
cannot catch this, because the off-topic answer may be perfectly true; only a direct
question-to-answer comparison does.

Two further signals were trialed for this family in an earlier revision. Both are
still computed on every query, but a later per-benchmark calibration retired them
from the composite, so they survive as recorded diagnostics rather than decision
signals; the composite stayed lean at nine.

\begin{description}[leftmargin=0pt, labelsep=0.5em, style=nextline]
  \item[Alias overlap.] Checks whether the named entities in the answer match the
  entities in the retrieved passages \emph{up to Wikidata aliases}, so that
  ``William Jefferson Clinton'' and ``Bill Clinton'' count as the same entity. It
  catches grounding failures that the entailment judge misses on surface form alone.
  \item[Entity-head consistency.] Reuses the three self-consistency chains and checks
  whether they name the \emph{same} head entity. It catches the case where the chains
  read as similar prose, so $\savg$ is high, yet they actually name different people
  or places. It costs no extra generation, since the chains already exist.
\end{description}

\begin{intuition}[Why they were trialed, and why the composite stayed at nine]
Each was introduced to close a specific, observed failure: a correct entity rejected
for using a different surface name, and three chains that looked consistent but
disagreed on who they were talking about. On the calibration fold, however, their
fitted contribution to the composite turned out negligible, so they were retired
from the trust score and kept only as cheap diagnostics. This is the panel
philosophy meeting its discipline: a specialist earns a vote in the composite only
if the calibration data show it adds discrimination, and these two did not, so the
deciding panel stays at nine.
\end{intuition}

\section{The confabulation early-exit}
\label{sec:verifier-earlyexit}

Stage 5 contains one hard rule that fires before any fusion: the confabulation
early-exit. If, at query time, the model's internal confidence is high while the
best passage barely supports the answer, the verifier short-circuits, returning a
discard verdict and sending the query to regeneration without computing the
composite at all.

\begin{figure}[t]
\centering
\begin{tikzpicture}[scale=5.4, font=\small\sffamily]
  % danger region: u_internal >= 0.70 AND p_ground_max <= 0.20
  \fill[cbPitfall!18] (0.70,0) rectangle (1.0,0.20);
  \fill[cbExample!8] (0,0.20) rectangle (1.0,1.0);
  \fill[cbExample!8] (0,0) rectangle (0.70,0.20);
  % boundaries
  \draw[cbPitfall!70!black, dashed, line width=0.7pt] (0.70,0) -- (0.70,0.20);
  \draw[cbPitfall!70!black, dashed, line width=0.7pt] (0.70,0.20) -- (1.0,0.20);
  % axes
  \draw[->, cbInk!65, line width=0.8pt] (0,0) -- (1.06,0) node[right, font=\footnotesize] {$\uint$};
  \draw[->, cbInk!65, line width=0.8pt] (0,0) -- (0,1.06) node[above, font=\footnotesize] {$\pgmax$};
  \foreach \x in {0.5,0.7,1.0} \draw[cbInk!40] (\x,0) -- (\x,-0.02) node[below,font=\scriptsize]{\x};
  \foreach \y in {0.2,0.5,1.0} \draw[cbInk!40] (0,\y) -- (-0.02,\y) node[left,font=\scriptsize]{\y};
  % labels
  \node[font=\footnotesize\bfseries, text=cbPitfall!70!black, align=center] at (0.85,0.10) {confabulation\\DISCARD + regenerate};
  \node[font=\footnotesize, text=cbExample!45!black, align=center] at (0.40,0.62) {proceed to the composite (\Cref{ch:composite})};
\end{tikzpicture}
\caption[The confabulation early-exit]{The one hard gate inside Stage 5. When
internal confidence is high ($\uint \ge 0.70$) but the best passage barely supports
the answer ($\pgmax \le 0.20$), the answer is a confident fabrication: the model is
sure of something the evidence does not back. The verifier discards it immediately
and regenerates, without bothering to compute the rest of the composite. Everywhere
else, it proceeds. (Schematic.)}
\label{fig:early-exit}
\end{figure}

\begin{precise}[Why this rule, and why only at query time]
The early-exit encodes the single most dangerous combination of signals: high
internal confidence paired with near-zero grounding is the textbook signature of
confabulation, an answer the model believes and the world does not support.
Catching it with a hard rule, rather than letting the calibrated composite weigh it,
guarantees that this specific pathology is always escalated rather than occasionally
averaged away. The rule fires only at query time, not during the cycle-boundary
re-verification of \Cref{ch:cycle}: after the model has been fine-tuned on its own
chains, high internal confidence is partly an artefact of having memorised them, so
the signal is no longer trustworthy as a confabulation flag, and re-verification
relies on its own grounding-based pruning instead.
\end{precise}

\section{What the verifier emits}
\label{sec:verifier-output}

The verifier returns a record carrying all the signals it computes, the nine that
feed the composite plus the recorded diagnostics, the canonical
answer it scored, the retrieved passages and atomic facts it used, and, after the
fusion of \Cref{ch:composite}, the composite trust score and the gate decision. Every
signal defaults to a neutral one-half on failure, so a missing dependency degrades
the panel gracefully rather than crashing it or silently scoring zero.

The two readings we have followed since \Cref{ch:bigpicture} diverge sharply at this
stage, and \Cref{fig:fingerprint} plots their reported signal values side by side.

\begin{figure}[t]
\centering
\begin{tikzpicture}
\begin{axis}[
  width=7.6cm, height=6.0cm,
  ybar, bar width=10pt,
  symbolic x coords={ground-max, s-avg, q-rel, composite},
  xtick=data, x tick label style={rotate=30, anchor=east, font=\scriptsize},
  ymin=0, ymax=1, ylabel={signal value}, font=\small\sffamily,
  tick label style={font=\scriptsize}, title={\footnotesize FEVER claim (deferred, later stored; illustrative)},
  ymajorgrids, grid style={cbInk!8},
  nodes near coords, nodes near coords style={font=\tiny, /pgf/number format/fixed, /pgf/number format/precision=2},
]
  \addplot[draw=cbExample!60!black, fill=cbExample!30] coordinates {
    (ground-max,0.94) (s-avg,0.85) (q-rel,0.69) (composite,0.51) };
  \draw[cbNumbers!70!black, dashed] (axis cs:ground-max,0.60) -- (axis cs:composite,0.60);
\end{axis}
\begin{axis}[
  xshift=8.4cm,
  width=7.6cm, height=6.0cm,
  ybar, bar width=10pt,
  symbolic x coords={u-token, s-avg, p-entail, ground-max, ground-atom, q-rel, composite},
  xtick=data, x tick label style={rotate=30, anchor=east, font=\scriptsize},
  ymin=0, ymax=1, font=\small\sffamily,
  tick label style={font=\scriptsize}, title={\footnotesize trivia question (abstained, deployed reading)},
  ymajorgrids, grid style={cbInk!8},
  nodes near coords, nodes near coords style={font=\tiny, /pgf/number format/fixed, /pgf/number format/precision=2},
]
  \addplot[draw=cbPitfall!70!black, fill=cbPitfall!25] coordinates {
    (u-token,0.02) (s-avg,0.80) (p-entail,0.12) (ground-max,0.11) (ground-atom,0.09) (q-rel,0.50) (composite,0.13) };
\end{axis}
\end{tikzpicture}
\caption[Two signal fingerprints]{Signal readings for the two running episodes. The
fact-checking claim (left) is an illustrative hand-authored walkthrough of the
defer-then-promote path: it grounds strongly in its best passage and is internally
consistent, lifting its composite into the deferral band, from which it is promoted to
storage on the next cycle. The trivia question (right) carries the actual deployed
cold-start verifier readings for the running episode, and is the chapter's own
``agreement is not support'' point made visible: the three chains agree with one
another ($\savg = 0.80$) yet do not entail the served answer ($\pentail = 0.12$), and
with grounding also near zero the calibrated composite settles at $\ustored = 0.13$,
below the deferral cut of $0.45$, and with $\pgmax = 0.11$ also below the grounding floor
of $0.20$, so the episode is abstained rather than silently discarded. The dashed line marks
the storage threshold (\Cref{ch:composite}).}
\label{fig:fingerprint}
\end{figure}

\section{Honest limitations}
\label{sec:verifier-retired}

Three caveats keep the verifier in proportion.

\begin{pitfall}[The verifier scores the answer, not the reasoning]
Every active signal scores the canonical \emph{answer} claim, never the individual
steps of the reasoning chain. The verifier therefore cannot distinguish a right
answer reached by sound reasoning from a right answer reached by a lucky wrong step,
nor can it credit a chain that reasons correctly but states the wrong final answer.
The model's logic is judged only insofar as it surfaces in the final answer text.
This is a known limitation, logged for future work; the panel is wide across
\emph{kinds} of evidence but shallow in the sense that it reads one claim, not a
derivation.
\end{pitfall}

\begin{pitfall}[Two signals were retired]
Two of the original signals no longer drive decisions. A semantic-entropy signal,
measuring the diversity of many resampled answers, was retired because its learned
weight in the composite shrank to nearly zero, so it added cost without information.
A contradiction probability was retired because the default short-hypothesis judge
reports only whether a passage \emph{supports} a claim, with no separate
contradiction verdict, so the signal was structurally always zero. Both remain in
the output record for diagnostics, but neither affects a decision; \caem{} relies on
\emph{low grounding} rather than an explicit contradiction veto.
\end{pitfall}

\begin{bythenumbers}[Verifier constants]
\begin{itemize}[leftmargin=1.3em, itemsep=1pt]
  \item Self-consistency chains: $M = 3$, sampled at temperature $0.7$, generated once and reused.
  \item Dropout passes: $K = 5$, dropout active, temperature $0.7$.
  \item Internal confidence: $\uint = 0.5\,\utoken + 0.5\,(1-\udrop)$.
  \item Grounding retrieval: top $20$ candidates, reranked to the best $3$ passages.
  \item Atomic decomposition: applied only to answers of at least $12$ tokens.
  \item NLI dispatch: MiniCheck for hypotheses $\le 408$ tokens, Frozen Qwen above.
  \item Early-exit: $\uint \ge 0.70$ and $\pgmax \le 0.20$ forces a discard at query time.
  \item Failure default: every signal returns a neutral $0.5$ on error.
\end{itemize}
\end{bythenumbers}

\begin{defence}[If an examiner asks: ``why so many signals, rather than one strong one?'']
Because the literature shows there is no single strong one: single-signal correctness
predictors barely beat chance (\Cref{ch:problem}). The contribution is not a better
individual detector but a \emph{panel} spanning four independent failure modes,
internal doubt, self-disagreement, ungrounded claims, and off-topic answers, so that
a wrong answer must defeat all of them at once. The diversity is the design. And the
panel is deliberately conservative at every turn: signals default to neutral on
failure, the entailment ensemble takes the minimum so all judges must agree, and a
hard early-exit rule guarantees the worst pathology is always escalated rather than
averaged away.
\end{defence}

The verifier has now produced a vector of calibrated-but-separate signals.
\Cref{ch:composite} takes up Stage 6: how those signals are each calibrated, fused
into a single trustworthy probability, and turned into the four-outcome decision that
governs storage.
